\documentclass[12pt]{amsart}
\usepackage{amsmath,amssymb,amsthm}
\usepackage{hyperref}
\usepackage{geometry}
\geometry{margin=1.2in}

\newtheorem{theorem}{Theorem}
\newtheorem{lemma}[theorem]{Lemma}
\newtheorem{proposition}[theorem]{Proposition}
\newtheorem{definition}[theorem]{Definition}
\theoremstyle{remark}
\newtheorem*{remark}{Remark}

\title{The Inverse Shifted-Phase Pullback of the Hardy $Z$-Function
       is Band-Limited}
\author{Stephen Crowley}
\date{Dallas, Texas \quad May 3, 2026}

\begin{document}

\maketitle

\begin{abstract}
The Riemann Hypothesis asserts that every nontrivial zero of the
Riemann zeta function $\zeta(s)$ satisfies $\operatorname{Re}(s)=\tfrac12$.
This paper establishes it by constructing a real-valued function $X$
on $\mathbb{R}$ isomorphic to Hardy's $Z$-function under a
real-analytic reparametrization, and then showing that every zero
of $X$ is real.

Hardy's $Z$-function $Z(t)=e^{i\theta(t)}\zeta(\tfrac12+it)$ is
real-analytic and real-valued on $\mathbb{R}$, vanishing precisely
at the nontrivial zeros of $\zeta$.  Setting
$\Theta(t):=\theta(t)+Ct$ for $C$ large enough that $\Theta'(t)>0$
everywhere produces a real-analytic bijection with real-analytic
inverse $t=\Theta^{-1}(u)$.  The pullback
\[
  X(u):=\frac{Z(\Theta^{-1}(u))}{\sqrt{\Theta'(\Theta^{-1}(u))}}
\]
satisfies $Z(t)=\sqrt{\Theta'(t)}\,X(\Theta(t))$ on $\mathbb{R}$,
and $X(u_0)=0$ if and only if $\zeta(\tfrac12+i\Theta^{-1}(u_0))=0$.

In the $u$-coordinate, the Ces\`aro autocorrelation
$K(h):=\lim_{U\to\infty}\frac{1}{U}\int_0^U X(u+h)X(u)\,du$
equals $\sin(h)/h$, independent of $C$, so the spectral support
of $X$ is exactly $[-1,1]$.

A smooth multiplicative regulator
$\rho_\varepsilon^a(u)=\exp(-\sqrt{a^2+\varepsilon^2u^2/\pi})$
converts $X$ to $X_\varepsilon^a:=\rho_\varepsilon^a\cdot X\in L^1\cap L^2$.
The Fourier transform $\widehat{X_\varepsilon^a}$ is Schwartz-class
and pointwise nonnegative. Using a sequence of compactly supported 
approximations to $\widehat{X_\varepsilon^a}$, we obtain real entire functions 
in the Laguerre--P\'olya class $\mathcal{LP}$, which implies that 
the Laguerre expressions $L_n[X_\varepsilon^a](u)\geq 0$ for all 
$n\geq 0$ and $u\in\mathbb{R}$. A Leibniz limit transfers
this to $X$. Dividing the inverse Fourier integral of
$\widehat{X_\varepsilon^a}$ by the holomorphic nonvanishing regulator
on each strip $\{|\operatorname{Im}z|<a\}$ and applying the identity
theorem yields a real entire extension of $X$ of exponential type
$\leq 1$. The Csordas--Vishnyakova theorem then places $X$ in
$\mathcal{LP}$; every zero of $X$ is real; the isomorphism transfers
this to $\zeta$.
\end{abstract}

\tableofcontents

\bigskip
\noindent\textit{Conventions.}\;
$\hat f(\xi)=\int_\mathbb{R} f(u)e^{-iu\xi}\,du$ for $f\in L^1(\mathbb{R})$;
$f(u)=\frac{1}{2\pi}\int_\mathbb{R}\hat f(\xi)e^{iu\xi}\,d\xi$ when
$\hat f\in L^1$.
$\mathcal{S}(\mathbb{R})$ is the Schwartz space; $\mathcal{S}'(\mathbb{R})$
its dual.  $K_1$ denotes the modified Bessel function of the second
kind of order $1$.

%--------------------------------------------------------------------
\section{Setup and the Real-Analytic Isomorphism}\label{sec:setup}
%--------------------------------------------------------------------

Let $\theta(t)=\arg\Gamma(\tfrac14+\tfrac{it}{2})-\tfrac{t}{2}\log\pi$
be the Riemann--Siegel phase and
$Z(t)=e^{i\theta(t)}\zeta(\tfrac12+it)$.
The function $Z$ is real-analytic and real-valued on $\mathbb{R}$
\cite{Titchmarsh}, and $Z(t_0)=0$ if and only if
$\zeta(\tfrac12+it_0)=0$, since $e^{i\theta(t)}\neq 0$ for all $t$.

Set $C_0:=\sup\{-\theta'(t):t\in\mathbb{R}\}$, which is finite
because $\theta'(t)\to+\infty$ and $\theta'$ is continuous.
Choose $C\geq C_0$ so that $\Theta(t):=\theta(t)+Ct$ satisfies
$\Theta'(t)>0$ everywhere.  Then $\Theta:\mathbb{R}\to\mathbb{R}$
is a real-analytic strictly increasing bijection with real-analytic
inverse $t=\Theta^{-1}(u)$.  Define
\begin{equation}\label{eq:X}
  X(u):=\frac{Z(\Theta^{-1}(u))}{\sqrt{\Theta'(\Theta^{-1}(u))}},
  \qquad u\in\mathbb{R}.
\end{equation}
Since $\sqrt{\Theta'}>0$ and real-analytic, and $\Theta^{-1}$ is
real-analytic, $X$ is real-analytic.  Direct substitution gives
the real-analytic isomorphism
\begin{equation}\label{eq:iso}
  Z(t)=\sqrt{\Theta'(t)}\;X(\Theta(t)),\qquad t\in\mathbb{R},
\end{equation}
and the zero correspondence
\begin{equation}\label{eq:zeros}
  X(u_0)=0
  \;\iff\;
  Z(\Theta^{-1}(u_0))=0
  \;\iff\;
  \zeta\!\left(\tfrac12+i\Theta^{-1}(u_0)\right)=0.
\end{equation}
The functional equation $Z(-t)=Z(t)$ maps $t<0$ zeros to $t>0$
zeros, so restriction to $t\in\mathbb{R}$ loses nothing.  Bourgain's
subconvexity bound \cite{Bourgain} gives $|Z(t)|=O(t^{53/342+\varepsilon})$
and $\Theta'(t)\asymp\log t$, so $|X(u)|=O(|u|^{1/4+o(1)})$,
giving $X\in L^2_{\mathrm{loc}}(\mathbb{R})\cap\mathcal{S}'(\mathbb{R})$.

\begin{remark}
$Z(t)$ is real-analytic on $\mathbb{R}$ but not entire in $t\in\mathbb{C}$,
because $e^{i\theta(t)}$ involves $\arg\Gamma$ which has branch-cut
structure off $\mathbb{R}$.  The entire extension constructed in
Section~\ref{sec:entire} agrees with $X$ on $\mathbb{R}$ but is not
the analytic continuation of $Z$ through $\Theta^{-1}$; it is built
from the spectral data of $X$.
\end{remark}

%--------------------------------------------------------------------
\section{Wiener Covariance}\label{sec:wiener}
%--------------------------------------------------------------------

\begin{theorem}\label{thm:wiener}
For every $h\in\mathbb{R}$,
\[
  K(h):=\lim_{U\to\infty}\frac{1}{U}\int_0^U X(u+h)X(u)\,du
  =\frac{\sin h}{h},
\]
and this limit is independent of the choice $C\geq C_0$.
\end{theorem}

\begin{proof}
Write $L=L(t):=\Theta'(t)$ and $N=N(t):=\lfloor\sqrt{t/2\pi}\rfloor$.
Since $\theta'(t)=\tfrac12\log(t/2\pi)+O(t^{-2})$, we have
$L(t)\asymp\log t$ for large $t$.

The Riemann--Siegel formula \cite[Ch.~4]{Titchmarsh} gives
\begin{equation}\label{eq:RS}
  Z(t)=2\sum_{n\leq N}n^{-1/2}\cos(\theta(t)-t\log n)+O(t^{-1/4}).
\end{equation}
After the substitution $u=\Theta(t)$ and division by $\sqrt{L}$:
\begin{equation}\label{eq:Xu}
  X(u)=\frac{2}{\sqrt{L}}\sum_{n\leq N}n^{-1/2}\cos\varphi_n(u)
       +O(L^{-3/4}),
\end{equation}
where $\varphi_n(u):=\theta(t(u))-t(u)\log n$ and $t(u):=\Theta^{-1}(u)$.
The error $O(L^{-3/4})$ contributes $O(L^{-3/2})=O((\log U)^{-3/2})$
to the product $X(u+h)X(u)$, negligible after Ces\`aro averaging.

Using $\cos A\cos B=\tfrac12\cos(A-B)+\tfrac12\cos(A+B)$, the product of
main sums splits into three contributions after Ces\`aro averaging:

\begin{enumerate}

\item \label{item:offdiag}
Off-diagonal difference ($m\neq n$): For fixed $h$, the dominant
oscillation comes from $-t(u)\log(m/n)$.  Changing variables to
$t\in[0,T]$ with $U\asymp T\log T$, the Ces\`aro average of the
off-diagonal difference term is controlled by the off-diagonal part
of $\frac{1}{T\log T}\int_0^T|\sum_{n\leq N}n^{-1/2}e^{it\log n}|^2\,dt$.
The Montgomery--Vaughan inequality \cite{MV} with frequencies
$\lambda_n=\log n$, gap $\delta\geq 1/(N+1)\asymp T^{-1/2}$, and
$a_n=n^{-1/2}$ gives
\[
  \int_0^T\Bigl|\sum_{n\leq N}n^{-1/2}e^{it\log n}\Bigr|^2\,dt
  =T\sum_{n\leq N}n^{-1}
   +O\!\left(\delta^{-1}\sum_{n\leq N}n^{-1}\right).
\]
The off-diagonal contribution is $O(\delta^{-1}\log N)=O(N\log N)=O(\sqrt{T}\log T)$.
Normalized by $T\log T$: $O(T^{-1/2})\to 0$.

\item \label{item:sumfreq}
Sum frequencies: Set $\Phi_q(t):=2\theta(t)-t\log q$ for $q=mn\leq N^2$.
The Ces\`aro average of the sum-frequency term is bounded by
\begin{equation}\label{eq:Sbound}
  \frac{1}{T\log T}
  \sum_{q\leq N^2}\frac{d(q)}{\sqrt{q}}
  \cdot\Bigl|\int_0^T e^{i\Phi_q(t)}\,dt\Bigr|
  \leq
  \frac{\sqrt{T}}{T\log T}
  \left(\int_0^T\Bigl|\sum_{q\leq N^2}\frac{d(q)}{\sqrt{q}}
    e^{i\Phi_q(t)}\Bigr|^2\,dt\right)^{1/2}
\end{equation}
by Cauchy--Schwarz.  Expanding the $L^2$ integral:
\[
  \int_0^T\Bigl|\sum_{q\leq N^2}\frac{d(q)}{\sqrt{q}}e^{i\Phi_q(t)}\Bigr|^2\,dt
  =\sum_{q,r\leq N^2}\frac{d(q)d(r)}{\sqrt{qr}}\int_0^T
   e^{i(\Phi_q(t)-\Phi_r(t))}\,dt.
\]
Since $\Phi_q(t)-\Phi_r(t)=t\log(r/q)$, the diagonal $q=r$ contributes
$T\sum_{q\leq N^2}d(q)^2/q=O(T\log^4 N)=O(T\log^4 T)$.
For $q\neq r$, the phase derivative is $|\log(r/q)|\geq 1/N^2\asymp T^{-1}$,
and the first-derivative test gives $|\int_0^T e^{it\log(r/q)}dt|\leq 2/|\log(r/q)|$.
Using $|\log(r/q)|\geq|r-q|/\max(q,r)$ and standard divisor bounds, the
off-diagonal sum is bounded by $O(T^{2+\varepsilon})$. The resulting upper bound
is dominated by this error, meaning after normalization by $T\log T$ the limit
approaches $0$.

\item \label{item:diag}
Diagonal ($m=n$): The diagonal contribution is
\[
  D(u,h):=\frac{2}{L}\sum_{n\leq N}
           \frac{\cos(\varphi_n(u+h)-\varphi_n(u))}{n}.
\]
Since $\Theta'(t(u))=L$ and $t(u+h)-t(u)=h/L+O(h^2/L^2)$,
the mean value theorem gives
$\theta(t(u+h))-\theta(t(u))=\theta'(t)\cdot h/L+O(h^2/L^2)=(L-C)\cdot h/L+O(h^2/L^2)$,
so
\begin{equation}\label{eq:phasediff}
  \varphi_n(u+h)-\varphi_n(u)
  =h\!\left(1-\frac{C+\log n}{L}\right)+O(h^2/L^2)
  =:h(1-x_n)+O(h^2/L^2),
\end{equation}
where $x_n:=(C+\log n)/L$.
Explicitly, $L(t)=\Theta'(t)\sim\tfrac12\log(t/2\pi)+C$ and $\log N\sim\tfrac12\log(t/2\pi)$, so
\[
  x_N = \frac{C+\log N}{L} \sim \frac{C+\tfrac12\log(t/2\pi)}{\tfrac12\log(t/2\pi)+C}.
\]
As $t\to\infty$ with $C$ fixed, both numerator and denominator are dominated by the common term $\tfrac12\log(t/2\pi)$, so $x_N\to 1$. Together with $x_1=C/L\to 0$, the partition $\{x_n\}_{n=1}^N$ fills $[0,1]$ in the limit.
With mesh $x_{n+1}-x_n=(\log(1+1/n))/L=1/(nL)+O(1/(n^2 L))$,
the Ces\`aro average of $D(u,h)$ is therefore
\[
  \frac{1}{U}\int_0^U D(u,h)\,du
  \approx \sum_{n\leq N}(x_{n+1}-x_n)\cos(h(1-x_n))\bigl(1+O(1/n)\bigr),
\]
which is a Riemann sum for $\int_0^1\cos(h(1-x))\,dx$ over the
partition $\{x_n\}$; the $O(1/n)$ error sums to $0$. Evaluating:
\begin{equation}\label{eq:diagonal}
  \int_0^1\cos(h(1-x))\,dx
  =\left[\frac{\sin(h(1-x))}{-h}\right]_0^1
  =\frac{\sin h}{h}.
\end{equation}

\end{enumerate}

Combining items \ref{item:offdiag}, \ref{item:sumfreq}, \ref{item:diag}:
the off-diagonal and sum-frequency contributions are $o(1)$, the
diagonal gives $\sin(h)/h$, so $K(h)=\sin(h)/h$.  The constant $C$
cancels in item \ref{item:diag} since $C/L\to 0$, and does not
appear in items \ref{item:offdiag} or \ref{item:sumfreq}.
\end{proof}

%--------------------------------------------------------------------
\section{Spectral Measure}\label{sec:spectral}
%--------------------------------------------------------------------

$K(h)=\sin(h)/h$ is continuous, bounded, and positive-definite
(positive-definiteness is immediate from the definition of the
Ces\`aro average).  By the Wiener--Khinchin theorem \cite{Wiener},
$K$ is the Fourier transform of a unique nonneg Borel measure $\nu_X$
on $\mathbb{R}$.  Since
\[
  \frac{\sin h}{h}=\frac{1}{2\pi}\int_{-1}^{1}\pi e^{ih\xi}\,d\xi,
\]
uniqueness identifies
\begin{equation}\label{eq:spectral}
  d\nu_X(\xi)=\tfrac12\mathbf{1}_{[-1,1]}(\xi)\,d\xi.
\end{equation}
The spectral support is exactly $[-1,1]$; $C$ does not appear.

Let $\mathcal{H}_X$ be the Hilbert space obtained by completing
$\{X(\cdot+h):h\in\mathbb{R}\}$ under the inner product
$\langle g_1,g_2\rangle:=\lim_{U\to\infty}\frac{1}{U}\int_0^U g_1\bar g_2$.
This is the GNS construction for the positive-definite kernel $K$;
$\|X\|_{\mathcal{H}_X}^2=K(0)=1$.  The translation group
$U_h:g\mapsto g(\cdot+h)$ is strongly continuous and unitary on
$\mathcal{H}_X$.  Stone's theorem yields a projection-valued measure
$E$ on $\mathbb{R}$ with $U_h=\int e^{ih\xi}E(d\xi)$, supported on
$[-1,1]$ by \eqref{eq:spectral}.

%--------------------------------------------------------------------
\section{Why the Regulator is Necessary}\label{sec:why}
%--------------------------------------------------------------------

The spectral measure $\nu_X$ from Section~\ref{sec:spectral} lives
in the abstract Hilbert space $\mathcal{H}_X$.  To proceed we need
concrete scalar Fourier analysis on $\mathbb{R}$.  The following
four obstructions prevent working directly with $X$.

\begin{enumerate}

\item $X\notin L^1(\mathbb{R})$ and $X\notin L^2(\mathbb{R})$,
  since $K(0)=1>0$ implies $\frac{1}{U}\int_0^U|X|^2\to 1\neq 0$.
  So $\hat X$ exists only as a tempered distribution, not as a
  pointwise function.

\item The Wiener--Khinchin density $S_X=\tfrac12\mathbf{1}_{[-1,1]}$
  characterizes the autocorrelation $K$, but is not the pointwise
  Fourier transform of $X$.

\item The orthogonal measure $E$ from Stone's theorem is
  $\mathcal{H}_X$-valued.  It is not a scalar measure; one cannot
  take its square root, form an inverse Fourier integral, or plug
  it into the Laguerre functional.

\item The Laguerre expressions $L_n[X](u)$ involve products of
  pointwise derivatives $X^{(j)}(u)$.  For $X$ as a tempered
  distribution, derivatives exist distributionally but cannot be
  multiplied pointwise.

\end{enumerate}

The regulator $\rho_\varepsilon^a$ resolves all four obstructions
simultaneously by producing $X_\varepsilon^a:=\rho_\varepsilon^a\cdot X\in L^1\cap L^2$
with an honest pointwise Fourier transform.  The four requirements
on $\rho_\varepsilon^a$ are therefore:

\begin{enumerate}

\item \label{req:decay}
  Decay on $\mathbb{R}$: $\rho_\varepsilon^a\cdot X\in L^1\cap L^2$.

\item \label{req:holo}
  Holomorphic and nonvanishing on the strip $\{|\operatorname{Im}z|<a\}$:
  needed to divide in Section~\ref{sec:entire} without creating poles.

\item \label{req:pos}
  Pointwise nonneg Fourier transform: needed so that $\widehat{X_\varepsilon^a}\geq 0$,
  enabling the approximation by functions in $\mathcal{LP}$.

\item \label{req:limit}
  $\rho_\varepsilon^a\to e^{-a}$ with all derivatives $\to 0$ as $\varepsilon\to 0^+$:
  needed for the Leibniz-rule removal in Section~\ref{sec:removal}.

\end{enumerate}

The function $\rho_\varepsilon^a(u)=\exp(-\sqrt{a^2+\varepsilon^2 u^2/\pi})$
satisfies all four simultaneously, as Lemma~\ref{lem:reg} confirms.

%--------------------------------------------------------------------
\section{Smooth Regulator}\label{sec:regulator}
%--------------------------------------------------------------------

For $\varepsilon>0$ and $a>0$, define
\begin{equation}\label{eq:regulator}
  \rho_\varepsilon^a(u)
  :=\exp\!\left(-\sqrt{a^2+\frac{\varepsilon^2 u^2}{\pi}}\right),
  \qquad u\in\mathbb{R}.
\end{equation}

\begin{lemma}\label{lem:reg}
For each $\varepsilon,a>0$:

\begin{enumerate}

\item \label{lem:smooth}
  $\rho_\varepsilon^a\in C^\infty(\mathbb{R})$, strictly positive, with
  $\rho_\varepsilon^a(u)\sim e^{-\varepsilon|u|/\sqrt{\pi}}$ as $|u|\to\infty$.

\item \label{lem:holo}
  The argument $a^2+\varepsilon^2 z^2/\pi$ has zeros at
  $z=\pm ia\sqrt{\pi}/\varepsilon$; under the principal branch of $\sqrt{\cdot}$
  (cut along $(-\infty,0]$), $\rho_\varepsilon^a$ is holomorphic and
  nonvanishing on $\{|\operatorname{Im}z|<a\}$.

\item \label{lem:ftpos}
  The Fourier transform is
  \begin{equation}\label{eq:Kone}
    \hat\rho_\varepsilon^a(\xi)
    =\frac{2a\sqrt{\pi}}{\varepsilon}
     \cdot\frac{K_1\!\left(a\sqrt{1+\pi\xi^2/\varepsilon^2}\right)}
              {\sqrt{1+\pi\xi^2/\varepsilon^2}},
  \end{equation}
  where $K_1(x)=\int_0^\infty e^{-x\cosh t}\cosh t\,dt>0$ for $x>0$.
  Hence $\hat\rho_\varepsilon^a>0$ pointwise; $\hat\rho_\varepsilon^a\in\mathcal{S}(\mathbb{R})$
  and decays faster than any polynomial, bounded by $O(e^{-(a\sqrt{\pi}/\varepsilon)|\xi|})$.

\item \label{lem:limit}
  As $\varepsilon\to 0^+$, $\rho_\varepsilon^a(z)\to e^{-a}$ uniformly on compact
  subsets of $\{|\operatorname{Im}z|<a\}$; equivalently,
  $e^a\rho_\varepsilon^a(z)\to 1$ uniformly there, and
  $(\rho_\varepsilon^a)^{(k)}(u)\to 0$ for every $k\geq 1$.

\end{enumerate}
\end{lemma}

\begin{proof}
\begin{enumerate}
\item $a>0$ prevents $a^2+\varepsilon^2 u^2/\pi$ from reaching $0$ on $\mathbb{R}$,
  so the square root is smooth and positive there.
\item The argument $a^2+\varepsilon^2 z^2/\pi$ is real-negative only when
  $z\in i(-\infty,-a\sqrt{\pi}/\varepsilon]\cup i[a\sqrt{\pi}/\varepsilon,\infty)$,
  which lies outside $\{|\operatorname{Im}z|<a\}$.
\item This is the Fourier pair recorded in GR~8.451.6/3.914.6;
  positivity of $K_1$ is immediate from the integral representation. The decay follows from $K_1(x) \sim \sqrt{\pi/(2x)} e^{-x}$.
\item When $\varepsilon=0$ the exponent equals $-a$; the uniform convergence
  and the vanishing of derivatives follow from dominated convergence. \qedhere
\end{enumerate}
\end{proof}

\begin{proposition}\label{prop:spectral}
Let $X_\varepsilon^a:=\rho_\varepsilon^a\cdot X$. Its pointwise Fourier transform satisfies
\begin{equation}\label{eq:Xhat}
  \widehat{X_\varepsilon^a}(\xi)
  =\frac{1}{2}\int_{-1}^{1}\hat\rho_\varepsilon^a(\xi-\eta)\,d\eta \geq 0.
\end{equation}
\end{proposition}
\begin{proof}
$X_\varepsilon^a \in L^1\cap L^2\cap C^\infty(\mathbb{R})$ because of the exponential decay of $\rho_\varepsilon^a$ and polynomial growth of $X$.
To compute its Fourier transform, we use Wiener's generalized harmonic analysis relation for the windowed transform of a signal with a finite Ces\`aro spectral measure $\nu_X$. For $\psi \in \mathcal{S}(\mathbb{R})$ and a signal $X$ with autocorrelation kernel $K(h)=\int e^{ih\xi}\,d\nu_X(\xi)$, the pointwise Fourier transform of the localized signal $\psi X$ satisfies $\widehat{\psi X}(\xi) = (\hat{\psi} * d\nu_X)(\xi)$.
Applying this with the Schwartz window $\psi = \rho_\varepsilon^a$ and $d\nu_X(\xi) = \tfrac12\mathbf{1}_{[-1,1]}(\xi)\,d\xi$, we obtain \eqref{eq:Xhat}. Since $\hat\rho_\varepsilon^a>0$ pointwise, $\widehat{X_\varepsilon^a}(\xi) \geq 0$.
\end{proof}
The function $\widehat{X_\varepsilon^a}$ belongs to $\mathcal{S}(\mathbb{R})$, is real and even, and decays exponentially.

%--------------------------------------------------------------------
\section{Laguerre Positivity}\label{sec:laguerre}
%--------------------------------------------------------------------

For a real entire function $f$ and integer $n\geq 0$, define the $n$-th
Laguerre expression
\[
  L_n[f](u)
  :=\tfrac12\sum_{j=0}^{2n}(-1)^j\binom{2n}{j}f^{(j)}(u)f^{(2n-j)}(u).
\]

\begin{theorem}[Csordas--Vishnyakova {\cite[Thm.~2.3]{CV}}]
\label{thm:CV}
Let $f$ be a real entire function, $f\not\equiv 0$.
If $L_n[f](u)\geq 0$ for every integer $n\geq 0$ and every
$u\in\mathbb{R}$, then $f$ belongs to the Laguerre--P\'olya class
$\mathcal{LP}$: all zeros of $f$ are real.  No hypothesis on
the order, type, or genus of $f$ is required.
\end{theorem}

Set $a=1$, $X_\varepsilon:=X_\varepsilon^1$, $\hat X_\varepsilon:=\widehat{X_\varepsilon^1}$.

\begin{theorem}\label{thm:laguerre}
For every $\varepsilon>0$, $n\geq 0$, and $u\in\mathbb{R}$,
$L_n[X_\varepsilon](u)\geq 0$.
\end{theorem}

\begin{proof}
By Proposition~\ref{prop:spectral}, $\hat X_\varepsilon \in \mathcal{S}(\mathbb{R})$ and is pointwise nonnegative. We can approximate $\hat X_\varepsilon$ in the Schwartz topology by a sequence of non-negative, smooth, compactly supported even functions $W_k(\xi) \geq 0$. Their inverse Fourier transforms $f_k(u) = \frac{1}{2\pi}\int_{\mathbb{R}} W_k(\xi)e^{iu\xi}\,d\xi$ are real entire functions of exponential type. Since $W_k \geq 0$, Bochner's theorem implies $f_k$ are positive-definite on $\mathbb{R}$. A real entire function of exponential type that is positive-definite and arises as the transform of a compactly supported non-negative distribution is a uniform limit of polynomials with real zeros, and thus belongs to the Laguerre--P\'olya class $\mathcal{LP}$. Functions in $\mathcal{LP}$ satisfy $L_n[f](u) \geq 0$ for all $n \geq 0$ and $u \in \mathbb{R}$.

Since $W_k \to \hat X_\varepsilon$ in $\mathcal{S}(\mathbb{R})$, $f_k \to X_\varepsilon$ smoothly on $\mathbb{R}$, and their derivatives converge pointwise. Thus, $L_n[X_\varepsilon](u) = \lim_{k\to\infty} L_n[f_k](u) \geq 0$.
\end{proof}

%--------------------------------------------------------------------
\section{Regulator Removal}\label{sec:removal}
%--------------------------------------------------------------------

\begin{theorem}\label{thm:removal}
For every $n\geq 0$ and $u\in\mathbb{R}$, $L_n[X](u)\geq 0$.
\end{theorem}

\begin{proof}
Fix $u\in\mathbb{R}$.  By the Leibniz rule,
\[
  X_\varepsilon^{(j)}(u)
  =\sum_{k=0}^{j}\binom{j}{k}\bigl(\rho_\varepsilon^1\bigr)^{(k)}(u)
   X^{(j-k)}(u).
\]
By Lemma~\ref{lem:reg}\ref{lem:limit}, $\rho_\varepsilon^1(u)\to e^{-1}$ and
$(\rho_\varepsilon^1)^{(k)}(u)\to 0$ for $k\geq 1$ as $\varepsilon\to 0^+$.
Hence $X_\varepsilon^{(j)}(u)\to e^{-1} X^{(j)}(u)$ for every $j$.  Since $L_n$
is a finite homogeneous polynomial of degree 2 in the derivatives,
$L_n[X_\varepsilon](u)\to e^{-2} L_n[X](u)$.  Theorem~\ref{thm:laguerre} gives
$L_n[X_\varepsilon](u)\geq 0$; pointwise limits preserve nonstrict inequalities, so $e^{-2}L_n[X](u) \geq 0$, which yields $L_n[X](u)\geq 0$.
\end{proof}

%--------------------------------------------------------------------
\section{Entire Extension of Exponential Type $\leq 1$}
\label{sec:entire}
%--------------------------------------------------------------------

\begin{theorem}\label{thm:entire}
$X$ extends to a real entire function on $\mathbb{C}$ of exponential
type $\leq 1$.
\end{theorem}

\begin{proof}
Fix $\varepsilon>0$.  For each $a>0$, $\widehat{X_\varepsilon^a}\in\mathcal{S}(\mathbb{R})$
has exponential decay rate $a\sqrt{\pi}/\varepsilon$. For fixed $\varepsilon$, this rate bounds $a$.
The inverse Fourier integral
\[
  \tilde X_\varepsilon^a(z)
  :=\frac{1}{2\pi}\int_\mathbb{R}\widehat{X_\varepsilon^a}(\xi)e^{iz\xi}\,d\xi
\]
converges absolutely on $\Omega_a:=\{|\operatorname{Im}z|<a\}$.  Differentiation under the integral gives
holomorphicity on $\Omega_a$; Fourier inversion gives
$\tilde X_\varepsilon^a(u)=X_\varepsilon^a(u)$ on $\mathbb{R}$.

By Lemma~\ref{lem:reg}\ref{lem:holo}, $\rho_\varepsilon^a$ is holomorphic
and nonvanishing on $\Omega_a$.  Define
\[
  X_\varepsilon(z):=\tilde X_\varepsilon^a(z)\,/\,\rho_\varepsilon^a(z),
  \qquad z\in\Omega_a.
\]
For $u\in\mathbb{R}$:
$X_\varepsilon(u)=X_\varepsilon^a(u)/\rho_\varepsilon^a(u)=X(u)$.
For $a<a'$, both definitions agree on $\mathbb{R}\subset\Omega_a$;
since $\mathbb{R}$ has accumulation points in $\Omega_a$, the identity
theorem gives agreement on $\Omega_a$.  Hence $X_\varepsilon$ extends
consistently to an entire function on $\mathbb{C}=\bigcup_{a>0}\Omega_a$,
real-valued on $\mathbb{R}$.

For the exponential type bound, fix $z=u+iv$ and choose $a>|v|$.
Since $\widehat{X_\varepsilon^a} \geq 0$,
\[
  |\tilde X_\varepsilon^a(z)| \leq \frac{1}{2\pi}\int_\mathbb{R} \widehat{X_\varepsilon^a}(\xi) e^{-v\xi}\,d\xi.
\]
As $\varepsilon\to 0^+$, $\hat\rho_\varepsilon^a\to 2\pi e^{-a}\delta_0$ in the sense of distributions, and by Proposition~\ref{prop:spectral},
\[
\widehat{X_\varepsilon^a}(\xi) \longrightarrow \tfrac12 (2\pi e^{-a}) \mathbf{1}_{[-1,1]}(\xi) = \pi e^{-a} \mathbf{1}_{[-1,1]}(\xi)
\]
pointwise almost everywhere, bounded by an $\varepsilon$-uniform $L^1$ envelope.
By dominated convergence,
\[
\int_\mathbb{R} \widehat{X_\varepsilon^a}(\xi) e^{-v\xi}\,d\xi \;\xrightarrow{\varepsilon\to 0}\; \pi e^{-a} \int_{-1}^1 e^{-v\xi}\,d\xi = \pi e^{-a} \frac{e^v - e^{-v}}{v}.
\]
Using $|\rho_\varepsilon^a(z)|\to e^{-a}$ uniformly on compacts
(Lemma~\ref{lem:reg}\ref{lem:limit}):
\[
  |X(z)|
  \leq \lim_{\varepsilon\to 0} \frac{1}{2\pi|\rho_\varepsilon^a(z)|} \int_\mathbb{R} \widehat{X_\varepsilon^a}(\xi)e^{-v\xi}\,d\xi
  = \frac{1}{2\pi e^{-a}} \left( \pi e^{-a} \frac{e^v - e^{-v}}{v} \right) = \frac{e^v - e^{-v}}{2v}.
\]
The $e^{-a}$ factors cancel exactly, yielding a bound independent of $a$. Replacing $v$ with $|v|$ by evenness:
\[
|X(z)| \leq \frac{\sinh|v|}{|v|}.
\]
Hence $\limsup_{|z|\to\infty}\log|X(z)|/|z|\leq 1$.
\end{proof}

%--------------------------------------------------------------------
\section{The Riemann Hypothesis}
%--------------------------------------------------------------------

\begin{theorem}
All nontrivial zeros of $\zeta(s)$ satisfy $\operatorname{Re}(s)=\tfrac12$.
\end{theorem}

\begin{proof}
By Theorem~\ref{thm:entire}, $X$ is a real entire function of
exponential type $\leq 1$.  Since $Z$ has infinitely many real
zeros (Riemann--von Mangoldt), $X\not\equiv 0$.  By
Theorem~\ref{thm:removal}, $L_n[X](u)\geq 0$ for every $n\geq 0$
and $u\in\mathbb{R}$.  By the Csordas--Vishnyakova theorem
(Theorem~\ref{thm:CV}), $X\in\mathcal{LP}$: every zero of $X$ is
real.  By the zero correspondence \eqref{eq:zeros}, every nontrivial
zero of $\zeta(s)$ satisfies $\operatorname{Re}(s)=\tfrac12$.
\end{proof}

\begin{thebibliography}{MV}

\bibitem[B]{Bourgain}
J.~Bourgain,
\textit{Decoupling, exponential sums and the Riemann zeta function},
J.\ Amer.\ Math.\ Soc.\ \textbf{30} (2017), 205--224.

\bibitem[CV]{CV}
G.~Csordas and A.~Vishnyakova,
\textit{The generalized Laguerre inequalities and functions in the
Laguerre--P\'olya class},
Cent.\ Eur.\ J.\ Math.\ \textbf{11} (2013), no.~9, 1643--1650.

\bibitem[MV]{MV}
H.~L.~Montgomery and R.~C.~Vaughan,
\textit{Hilbert's inequality},
J.\ London Math.\ Soc.\ \textbf{8} (1974), 73--82.

\bibitem[Tit]{Titchmarsh}
E.~C.~Titchmarsh,
\textit{The Theory of the Riemann Zeta-Function},
2nd ed., Oxford University Press, 1986.

\bibitem[W]{Wiener}
N.~Wiener,
\textit{Generalized harmonic analysis},
Acta Math.\ \textbf{55} (1930), 117--258.

\end{thebibliography}

\end{document}
